\subsection*{The binary cover, periodic braids, and the normalizer test}

\statement{B.8}{Proved}{Explicit periodic braids over the icosahedral action}
Put \(Z=\Delta_{12}^{2}\), where
\(\Delta_{12}=(\sigma_1\cdots\sigma_{11})(\sigma_1\cdots\sigma_{10})
\cdots(\sigma_1\sigma_2)\sigma_1\) is the Garside half twist in the
discrete Artin braid group \(B_{12}\). There are explicit braids
\(\beta_2=\Delta_{12}\) and \(\beta_3\) satisfying
\[
 \beta_2^2=\beta_3^3=Z.
\]
We specify \(\beta_3\) by listing the indices of its positive letters,
in multiplication order:
\[
\begin{gathered}
1,11,5,7,10,2,9,3,8,4,2;\\
7,3,6,2,5,4,3,2,7,10,6;\\
1,7,9,10,8,5,9,7,2,6,10;\\
5,4,7,3,5,8,4,2,5,7,6.
\end{gathered}
\]
With \(a=\beta_2\beta_3^{-1}\) and
\[
 R_{12}=\sigma_{11}\sigma_{10}\cdots\sigma_2\sigma_1^2
              \sigma_2\cdots\sigma_{10}\sigma_{11},
\]
the stronger Artin-group identity is
\[
 a^5=Z R_{12}^{-1}.
\]
Consequently the images in \(B_{12}(S^2)\), where \(R_{12}=1\),
satisfy \(a^5=\beta_3^3=(a\beta_3)^2=Z\) and generate the binary
icosahedral group of order \(120\). Their mapping-class image is
the icosahedral \(A_5\).
\begin{proof}
The words were discovered by following rigid rotations of an icosahedron
about an edge axis and an adjacent face axis. No numerical approximation
is used in their certification. The program
\texttt{check\_icosahedral\_braids.py} uses Artin's faithful action on
the discrete free group on \(x_1,\ldots,x_{12}\):
\[
 \sigma_i(x_i)=x_ix_{i+1}x_i^{-1},\quad
 \sigma_i(x_{i+1})=x_i,\quad
 \sigma_i(x_j)=x_j\quad(j\ne i,i+1).
\]
Free reduction verifies all three displayed identities, including the
identity of the discovered 66-letter word with \(\Delta_{12}\).
The code and its complete output are included. After imposing
\(x_1\cdots x_{12}=1\), it also verifies directly that \(a^5\)
is inner, with conjugating word \(x_1\cdots x_{11}\).

The permutations of \(\beta_2,\beta_3,a\) have orders \(2,3,5\),
and exhaustive permutation multiplication gives a group of order
sixty. The triangle presentation of B.4 therefore identifies the
mapping-class image with \(A_5\), not with a proper quotient.
The standard sphere-braid exact sequence
\[
 1\longrightarrow\langle Z\rangle\cong C_2
 \longrightarrow B_{12}(S^2)
 \longrightarrow\Gamma_{0,[12]}\longrightarrow1
\]
is recalled and derived geometrically in
\citet[Sections 1--3]{GoncalvesGuaschi2008}.
Our subgroup contains its nonidentity kernel element \(\beta_2^2=Z\),
so its order is \(120\). Independently,
\texttt{check\_icosahedral\_groups.py} verifies the order of
\(\langle a,b\mid a^5=b^3=(ab)^2\rangle\) by Todd--Coxeter
enumeration and realizes it as \(\SL_2(\mathbb F_5)\).

For the geometric identifications, the orbit of the unordered vertex
set is \(SO(3)/A_5\); its universal cover is \(SU(2)\), with deck
group the inverse image of \(A_5\), namely the binary icosahedral
group. It is the Poincar\'e sphere. In the Birman extension with one
additional distinguished puncture, the inverse image of \(A_5\)
is the orbifold fundamental group of the quotient of the punctured
sphere, of signature \((0;2,3;\infty)\). It is \(C_2*C_3\).
The preimage in the Artin group is its central extension by the
boundary twist. The two primitive cone lifts obey
\(\beta_2^2=\beta_3^3=Z\), so its presentation is
\(\langle\beta_2,\beta_3\mid\beta_2^2=\beta_3^3\rangle\cong B_3\).
To check that these generators give the entire preimage, first
work modulo \(Z\). In the orbifold group \(C_2*C_3\), the
point-pushing meridian \(R_{12}^{-1}\) is conjugate to \(\tau^5\).
The identity \(a^5=R_{12}^{-1}\), cyclic centralizers of
parabolics, and unique roots in a parabolic cyclic subgroup make
\(a\) a conjugate of the primitive cusp generator \(\tau\).
In the Bass--Serre tree of \(C_2*C_3\), let \(v_2,v_3\) be
the unique vertices fixed by the two elliptics. Their types differ,
so the vertices are distinct. Since edge stabilizers are trivial,
neither elliptic fixes the edge incident to its endpoint on
\([v_2,v_3]\). Concatenating this segment with its translates
therefore gives an axis for their product, with translation
length \(2d(v_2,v_3)\). The primitive cusp \(\tau\) has
translation length two, and so does its conjugate \(a\).
Thus \(v_2,v_3\) are adjacent. The ambient group is transitive
on edges; their stabilizers are consequently a simultaneous
conjugate of the two factors at the fundamental edge.
Each elliptic generates its entire vertex stabilizer, so the
two elliptics generate all of \(C_2*C_3\).
Since the Artin subgroup also contains \(Z\), it is the entire
central preimage. Removing the order-five
Seifert fibre gives the trefoil exterior, the Seifert manifold over
a disk with cone orders two and three. These descriptions agree
with \citet[Descriptions 6--8]{KirbyScharlemann1979}.
\end{proof}

\statement{B.9}{Proved}{The normalizer formulation and its three-coset test}
Let \(E=\widehat\Gamma_{0,[12],1}\), with the usual completed
Birman homomorphism \(p:E\to\widehat\Gamma_{0,[12]}\), and identify
\(\widehat B=p^{-1}(A_5)\). Restriction of conjugation gives
\[
 N_E(\widehat B)\longrightarrow\Out(\widehat B).
\]
Its image is exactly \(\mathcal K\). For a multiplier-one
representative normalized by \(\Phi(\tau)=\tau\), with frames
\(\Phi(s)=\beta_\Phi s\beta_\Phi^{-1}\) and
\(\Phi(r)=A_\Phi r A_\Phi^{-1}\), membership in \(\mathcal K\)
is equivalent to
\[
 C_E(\tau)\cap\beta_\Phi C_E(s)\cap A_\Phi C_E(r)\ne\varnothing.
\]
\begin{proof}
Genus-zero congruence and the completed point-pushing sequence identify
the kernel of \(p\) with \(\widehat K\) and identify
\(p^{-1}(A_5)\) with the completion of the discrete orbifold
extension. Asada's Theorems 1, 2 and 3A prove the necessary
faithfulness and the induction under addition of a puncture
\citep{Asada2001}. A lift of a geometric outer automorphism of
\(\widehat K\) is therefore an element of \(E\).

If that outer automorphism centralizes \(A_5\), the extension
comparison B.5 shows that its lift normalizes \(\widehat B\)
and restricts, modulo inner automorphisms, to the corresponding
element of \(\mathcal H\). Conversely a normalizing element gives
a geometric outer automorphism of \(\widehat K\) whose induced
automorphism of \(A_5\) is inner: its multiplier is one, and B.6
applies. Thus its restriction is in \(\mathcal K\).

If \(\Phi\) is induced by such conjugation only up to an inner
automorphism of \(\widehat B\), multiply the normalizing element
by the inverse inner conjugator. We then seek \(u\in E\) satisfying
\(u\tau u^{-1}=\tau\),
\(usu^{-1}=\beta_\Phi s\beta_\Phi^{-1}\), and
\(uru^{-1}=A_\Phi r A_\Phi^{-1}\).
These are precisely the three coset conditions. Conversely the
conditions give agreement on the topological generators \(s,r\),
so give the required automorphism on all of \(\widehat B\).
\end{proof}

\statement{B.10}{Proved}{The properly typed tripod tests}
For \(\Phi\in\mathcal H\) and a fixed three-subset \(t\) of the
twelve labelled cusps, let \(N_t\trianglelefteq\widehat K\) be
the closed normal subgroup generated by all other cusp subgroups.
Then restriction and quotient define a continuous homomorphism
\[
 \mathfrak T_t:\mathcal H\longrightarrow
 \Out_{\partial}(\widehat K/N_t)\cong\mathcal P.
\]
The last identification requires an ordered identification of the
three surviving cusps with \(0,1,\infty\). We have
\(\mathfrak T_t(\mathcal K)=1\) and
\(\mathcal K\cap G_{\mathbb Q}=1\), for the specified faithful
Galois embedding. The same test is defined on any subgroup fixing
the three-subset setwise; it is not automatically a map for every
element of the larger group in B.6 and every labelled three-subset.
\begin{proof}
Purity in B.7 makes \(N_t\) invariant. A quotient of a free
profinite punctured-sphere group by the discarded peripheral normal
closure is the full free profinite group of the retained punctured
sphere; this follows directly from its presentation. Inner changes
of representatives descend to inner changes, so the map is well
defined and continuous. A pure geometric mapping class becomes the
identity after forgetting down to three punctures. Continuity gives
the assertion for its completion and hence for \(\mathcal K\).

All twelve cusps of the chosen modular model are rational over
\(\mathbb Q(\zeta_5)\). For an arithmetic element fixing that
field, quotienting is the functorial outer Galois action on the
three-punctured sphere. A fractional linear transformation defined
over that field identifies the three chosen rational cusps with
\(0,1,\infty\); thus the test is the usual Galois outer action,
with the stated coordinate choice. Belyi faithfulness makes its
kernel trivial. A Galois element in \(\mathcal K\) has multiplier
one and therefore belongs to \(G_{\mathbb Q(\zeta_5)}\), proving
the intersection assertion.

For \(\lambda\equiv-1\pmod5\), the equation in B.7 permits the
nontrivial normalizer action on a five-cycle. It can exchange
antipodal cusp labels. An arbitrary three-subset then need not be
preserved, so the unqualified version of P3 has no such quotient
endomorphism as its target. One instead obtains a map between
different labelled quotients, or restricts to their stabilizer.
\end{proof}

\statement{B.11}{Proved}{The topological quintic resolvent and its distinct quotient types}
The preimage \(\Gamma_T\leq B\) of an \(A_4\leq A_5\) has
signature \((0;2,3,3;\infty)\) and presentation
\(C_2*C_3*C_3\). Its degree-five modular map has the model
\[
 j=t^3(t^2+5t+40),\qquad
 j-1728=(t-3)(t^2+4t+24)^2.
\]
The two unramified points above \(j=0\), hence the two order-three
orbifold points, are conjugate over \(\mathbb Q(\sqrt{-15})\).
Filling either order-three point gives \(C_2*C_3\); filling the
order-two point gives \(C_3*C_3\), not \(C_2*C_3\).
\begin{proof}
The degree-five permutation action on \(A_5/A_4\) gives cycle
types \(2^2 1\), \(3\,1^2\), and \(5\) at the three base
points. The orbifold cover consequently has one cone point of
order two, two of order three, and one cusp. The orbifold Euler
characteristic, or its standard punctured presentation, gives the
stated free product. The polynomial identity is verified by
expansion. The roots of \(t^2+5t+40\) have discriminant
\(25-160=-135\), hence the specified quadratic field.
Killing a cone free factor in a free product leaves the other
factors, both discretely and after profinite completion by its
universal property. Their types distinguish the three targets.
\end{proof}

\statement{B.12}{Open}{The proposed identities on the full Grothendieck--Teichm\"uller group}
B.53 below proves P4's exact quadratic-character domain, both
quintic homomorphisms, their arithmetic identity, and the kernel
inclusion for $\mathcal K$. Its third map has the corrected
weighted-orbifold target $\operatorname{Out}_{\partial}(C_3\amalg C_3)$.
B.55 constructs a nonempty set of full R2 lifts $\Phi\in\mathcal H$
for which $\mathfrak T_\pm(\Phi)\ne\Phi$, proving the first
assertion of P5. B.57 refutes P7's full common-kernel equality
by a separate full central-extension construction.
What remains undecided here is the assertion that
$\mathfrak T_\pm$ fixes every element of $\widehat{GT}\cap D$.
That assertion does not follow from the arithmetic restriction
or from a linearized calculation.

The exact source distinction is decisive: Nakamura--Schneps,
Theorem 1.1 and Definition 1.1, prove their additional relations
on the Galois image and use them to define the subgroup
\(I\Gamma\); Theorem 1.2 proves that it is a subgroup containing
that image \citep{NakamuraSchneps2000}. These statements do not
assert that all of \(\GT\) satisfies the new relations.
Hoshi--Mochizuki--Tsujimura, Definition 3.3, retain the additional
COF and RGC conditions in their BGT framework
\citep{HoshiMochizukiTsujimura2025}. Replacing those groups by
\(\GT\) would remove a hypothesis. The unresolved requirement
is a proof of that replacement, or a concrete element violating
the claimed identity, together with its typed restriction maps.

\statement{B.13}{Proved}{Rigid cyclic subconfigurations impose genuine necessary equations}
Let \(Q\) be a subset of the twelve vertices invariant under a
cyclic subgroup \(H\leq A_5\), such that the quotient of the
\(Q\)-punctured sphere by \(H\), including its remaining cone
points, has three distinguished points. Forgetting the other
punctures maps \(\mathcal K\) to the identity in
\(\widehat\Gamma_{0,[Q]}\). This includes a pentagon, a pentagon
with one fixed pole, and a pentagon with both fixed poles.
\begin{proof}
Because \(\mathcal K\) is pure, forgetting defines a continuous
homomorphism on it and on the stabilizer of \(Q\). It is
\(H\)-equivariant by functoriality, so the image centralizes \(H\).
The discrete normalizer of this cyclic action modulo deck
transformations injects into the mapping class group of the
three-distinguished-point quotient, a finite group. Thus the
discrete centralizer is finite. The single prime-power case of
Boggi--Zalesskii, Theorem A, gives its completed centralizer as
that finite discrete group. Its intersection with the pure
mapping class group is trivial, since a finite-order mapping
class of a punctured sphere fixing at least three punctures is
the identity. Purity descends under forgetting, so the image is
trivial. No nonsolvable centralizer comparison occurs in this
argument.
\end{proof}

\statement{B.14}{Proved}{The presentation complex and the limitation of its arithmetic action}
The presentation complex of
\(\langle a,b\mid a^5b^{-3},\ b^3(ab)^{-2}\rangle\)
is homotopy equivalent to the Poincar\'e homology sphere with
an open ball removed. In a trefoil presentation
\(\langle b,d\mid b^3=d^2\rangle\), put
\(\mu=db^{-1}\) and \(z=b^3=d^2\).
Then \(\mu\) is a meridian, and the additional two-cell is
attached along \(\mu^5z^{-1}\).
\begin{proof}
The substitution \(d=ab\), \(a=db^{-1}\) is a Nielsen
change of generators and converts the two relators into
\(\mu^5z^{-1}\) and \(b^3d^{-2}\), with the usual inversion
or multiplication of relators giving elementary changes of the
presentation complex. The one-relator complex for the trefoil is
homotopy equivalent to a spine of its exterior: decompose it into the two
solid-torus pieces of the torus knot and retain their annular
intersection. To identify the peripheral slope, use Wirtinger
generators \(x,y\) with \(xyx=yxy\), and put
\(b=xy\), \(d=xyx\). Then \(db^{-1}=bxb^{-1}\), explicitly
a conjugate of a meridian. The regular boundary fibre has
torus-framing linking number six, so the preferred longitude is
\(z\mu^{-6}\). Thus \(\mu^5z^{-1}\) is the negative
of the meridian-plus-longitude filling slope.
Attaching its two-handle and omitting the final three-handle
gives precisely the filled manifold minus a ball. Homotopy
invariance of attaching a cell gives the asserted equivalence
with the presentation complex. The trefoil surgery description in
\citet[Descriptions 6--8]{KirbyScharlemann1979} identifies the
filling with the Poincar\'e sphere, with orientation depending
on the handedness conventions. The homotopy assertion is
orientation independent.
\end{proof}

\statement{B.15}{Proved}{Triviality on the binary quotient means triviality of an outer action}
For a multiplier satisfying \(\lambda\equiv\pm1\pmod5\),
the induced outer automorphism on the binary icosahedral quotient
is trivial. A chosen representative that acts identically on
\(A_5\) acts identically on \(\SL_2(\mathbb F_5)\).
An arbitrary automorphism representative need not be the identity.
\begin{proof}
This is the functorial action on the central double cover, not
a quotient of \(\widehat B\) itself. Every automorphism of
\(A_5\) lifts: identify \(\Aut(A_5)\) with
\(\operatorname{PGL}_2(\mathbb F_5)\), and lift by conjugation
in \(\operatorname{GL}_2(\mathbb F_5)\).
Its central kernel has order two, and every automorphism fixes
that kernel pointwise. An automorphism inducing the identity on
\(A_5\) has the form \(g\mapsto\epsilon(g)g\), where
\(\epsilon:\SL_2(\mathbb F_5)\to C_2\) is a homomorphism.
The binary group is perfect: the exponent-sum matrix of the
displayed presentation is
\(\left(\begin{smallmatrix}5&-3\\-2&1\end{smallmatrix}\right)\),
whose determinant is \(-1\). Hence \(\epsilon=1\).
The action on \(A_5\) is inner by B.6; adjusting by a lift
of that inner automorphism gives the stated identity representative.
Without adjustment, nontrivial inner conjugations are explicit
counterexamples to literal identity on every element. This is
why the outer/genuine distinction cannot be dropped from B.8
of the prompt.
\end{proof}
